\documentclass[12pt]{article}
\usepackage[utf8]{inputenc}
\usepackage[T1]{fontenc}
\usepackage[margin=1in]{geometry}
\usepackage{setspace}
\usepackage{xcolor}
\usepackage[colorlinks=true, urlcolor=blue!60!black, linkcolor=blue!60!black]{hyperref}
\setlength{\parindent}{0em}
\setlength{\parskip}{0.45em}
\renewcommand{\baselinestretch}{1.0}
\begin{document}
\begin{center}
{\Large \textit{DeShaney v. Winnebago County Department of Social Services}, 489 U.S. 189 (1989)\par}
\vspace{0.3em}
{\normalsize Case Brief --- One-Page Edition\par}
\end{center}

\textbf{Memory Jogger.} ``Poor Joshua'': no constitutional duty to protect against private violence (\href{https://supreme.justia.com/cases/federal/us/489/189/}{489 U.S. 189}).

\textbf{Facts.} County social workers documented repeated signs that Randy DeShaney was abusing his son Joshua, briefly took custody, returned him, and then did nothing; Randy beat Joshua into permanent brain damage.

\textbf{Procedural History.} Summary judgment for the county on the \S~1983 claim; the Seventh Circuit affirmed.

\textbf{Votes.} Affirmed, 6--3: Rehnquist for the Court; Brennan and Blackmun dissented.

\textbf{Holding.} The Due Process Clause limits state action and creates no affirmative right to protection, absent custody or a danger the State created or worsened.

\textbf{Reasoning.} The Clause protects ``the people from the State, not to ensure that the State protected them from each other.'' Custody cases like \textit{Estelle v. Gamble} and \textit{Youngberg v. Romeo} did not apply because Joshua was in his father's custody, and returning him left him no worse off. The Amendment ``does not transform every tort committed by a state actor into a constitutional violation.'' Brennan stressed that the State monopolized rescue; Blackmun decried ``sterile formalism.''

\textbf{Examples.} \emph{Same-Side:} slow police response to a protective-order violation. \emph{Opposite-Side:} a trooper stranding a child on a highway. \emph{Fence-Sitter:} abuse in state foster care.

\textbf{Critique.} Echoing Brennan, critics call ``inaction'' a baseline choice where the State gave Randy custody and monopolized rescue; the ruling also sits uneasily with the Amendment's origins in protecting Black citizens from Klan violence.

\textbf{Key Quotation.} ``While the State may have been aware of the dangers that Joshua faced in the free world, it played no part in their creation, nor did it do anything to render him any more vulnerable to them.'' --- Chief Justice Rehnquist, for the Court, 489 U.S. at 201.

\end{document}
